\newgeometry{left=18mm,right=18mm, top=16mm, bottom=16mm}
\thispagestyle{empty}
\section*{Zusammenfassung}

%new version
Die Nukleosynthese schwerer Elemente (A > 56) jenseits von Eisen zählt zu den faszinierendsten und zugleich komplexesten Prozessen der nuklearen Astrophysik. Der schnelle Neutroneneinfangprozess (r-Prozess), der für die Erzeugung des Großteils der im Universum beobachteten schweren Elemente verantwortlich ist, findet in astrophysikalischen Umgebungen mit extrem hohen Neutronenflussdichten statt, wie etwa bei Neutronensternverschmelzungen oder Kernkollaps-Supernovae. Die Untersuchung der Nuklearstruktur exotischer, neutronenreicher Kerne weit abseits der Stabilitätslinie ist entscheidend für das Verständnis des r-Prozess-Pfades und liefert wesentliche Informationen zur Einschränkung der Zustandsgleichung (Equation of State, EoS) von Neutronensternen – astrophysikalischen Objekten von herausragender Bedeutung. Die EoS bestimmt die makroskopischen Eigenschaften von Neutronensternen, einschließlich ihrer Massen-Radius-Relation. Während astronomische Beobachtungen makroskopische Einschränkungen für Neutronensternmodelle liefern, stellt die experimentelle Kernphysik unverzichtbare Informationen über die mikroskopischen nuklearen Parameter bereit, auf die diese Modelle angewiesen sind.\\
Eine umfassende Beschreibung des r-Prozesses erfordert zudem detaillierte Kenntnisse der Nuklearstruktur und der zugrunde liegenden Nukleon-Nukleon-Wechselwirkungen. Insbesondere eine präzise Modellierung der Spaltprozesse, die am Ende des r-Prozess-Pfades stattfinden und über Spaltrecycling erneut Material in den r-Prozess einspeisen, ist von herausragender Bedeutung, um die beobachteten Häufigkeiten schwerer Elemente konsistent erklären zu können.\\
Diese Arbeit mit dem Titel \textit{„Nuclear Structure Investigations of Light Nuclei with the R3B Experiment“} präsentiert detaillierte Untersuchungen moderner experimenteller Techniken zur hochpräzisen Bestimmung zentraler nuklearer Eigenschaften. Die Experimente wurden mit dem R3B-Setup (Reactions with Relativistic Radioactive Beams) im Rahmen der Kommissionierungsexperimente S444 und S455 durchgeführt, die Teil der FAIR Phase-0 Kampagne am GSI waren.\\
Zusammenfassend präsentiert diese Arbeit eine umfassende Analyse nuklearer Streuobservablen und Reaktionsmechanismen innerhalb des R$^3$B-Setups. Sie umfasst (i) hochpräzise Messungen der Ladungs\-wechsel- und Totalwechselwirkungsquerschnitte für $^{12}$C + $^{12}$C-Kollisionen im Energiebereich von 400–800 MeV/Nukleon; (ii) die Validierung der Quasi-freien Streureaktion $^{12}$C(p,2p)$^{11}$B in inverser Kinematik als Werkzeug zur Untersuchung der Einteilchenstruktur von Kernen, einschließlich der ersten eindeutigen Rekonstruktion des Gammaspektrums von  $^{11}$B unter Verwendung des CALIFA-Kalorimeters; (iii) eine Analyse der Spaltung mittels quasi-freier Streuung als neuartigen Ansatz zur umfassenden Charakterisierung des Spaltprozesses, komplementär zur Coulomb-exzitationsinduzierten Spaltung (Coulex-Fission); sowie (iv) die Entwicklung fortgeschrittener Cluster-Rekonstruktionsverfahren für den CALIFA-Kalorimeter mithilfe von Methoden des maschinellen Lernens.\\
In ihrer Gesamtheit tragen diese Untersuchungen zur Verfeinerung experimenteller Methoden zur Erforschung exotischer Kerne bei. Diese Entwicklungen sind wegweisend für die "Early Science" Kampagne, die 2028 bei FAIR beginnt.

